\chapter{آرایه‌ها}
\label{ch:arrays}

\begin{goals}
\begin{itemize}
  \item نگه‌داری چند مقدار هم‌نوع در یک متغیر
  \item دسترسی به هر عضو با شماره (اندیس) و تغییر دادن آن
  \item پیمایش آرایه با \fcode{هر}
  \item الگوهای پایه: جمع، میانگین، بیشترین، شمارش و جست‌وجو
  \item کار با چند آرایه‌ی هم‌راستا
\end{itemize}
\end{goals}

\section{سی متغیر یا یک آرایه؟}

فرض کنید می‌خواهید نمره‌های سی دانش‌آموز را نگه دارید و میانگین بگیرید.
با آنچه تا حالا داریم، باید سی متغیر بسازید: \fcode{نمره۱}، \fcode{نمره۲}،
\ldots{} و سی بار آن‌ها را جمع کنید. و اگر کلاس سی‌ویک نفر شد؟ روشن است که این
راه به جایی نمی‌رسد.

\textbf{آرایه} پاسخ این مشکل است: یک متغیر که به جای یک مقدار، فهرستی از
مقدارهای هم‌نوع را نگه می‌دارد. به جای سی جعبه‌ی جداگانه، یک قفسه با سی
خانه‌ی شماره‌دار.

\section{ساختن آرایه و دسترسی به عضوها}

آرایه را با نوشتن مقدارها میان دو کروشه و با ویرگول می‌سازیم:

\program{ساختن آرایه}
\begin{salam}
تابع اصلی:
    نمرات := [۱۷, ۱۵, ۱۹, ۱۲, ۲۰]
    چاپ نمرات
    چاپ "اولین نمره:", نمرات[۰]
    چاپ "سومین نمره:", نمرات[۲]
    چاپ "تعداد نمره‌ها:", نمرات.طول()
تمام
\end{salam}

\begin{outputfa}
[17, 15, 19, 12, 20]
اولین نمره: 17
سومین نمره: 19
تعداد نمره‌ها: 5
\end{outputfa}

چند چیز در این برنامه تازه است:
\begin{itemize}
  \item \fcode{چاپ نمرات} کل آرایه را با همان شکل کروشه‌ای چاپ می‌کند.
        برای نگاه سریع به محتوای آرایه بسیار مفید است.
  \item \fcode{نمرات[۰]} عضو اول آرایه است. عدد داخل کروشه \textbf{اندیس}
        نام دارد و \emph{از صفر} شروع می‌شود: عضو اول اندیس ۰ دارد، عضو
        دوم اندیس ۱ و همین‌طور تا آخر. پس \fcode{نمرات[۲]} عضو سوم است.
  \item \fcode{نمرات.طول()} تعداد عضوهای آرایه را می‌دهد. با همان نقطه‌ای
        نوشته می‌شود که در فصل \ref{ch:types} برای \fcode{به صحیح} دیدیم.
\end{itemize}

\begin{note}[چرا از صفر؟]
اندیس را «فاصله از آغاز» بدانید نه «شماره‌ی ردیف». عضو اول، صفر خانه از
آغاز فاصله دارد. با این نگاه، شروع از صفر طبیعی است و در فصل
\ref{ch:loops} هم دیدید که شمارنده‌ی \fcode{تکرار} از صفر شروع می‌شود؛
این دو با هم جور درمی‌آیند. آخرین عضو آرایه‌ای با ۵ عضو، اندیس ۴ دارد، یعنی
همیشه \fcode{طول() - ۱}.
\end{note}

\begin{warn}
اگر اندیسی بیرون از محدوده بدهید (مثلاً \fcode{نمرات[۵]} برای آرایه‌ای
با ۵ عضو، یا یک اندیس منفی)، برنامه هنگام اجرا با خطا متوقف می‌شود. این
یکی از رایج‌ترین خطاهای کار با آرایه است و بیشتر وقت‌ها از فراموش کردن
«از صفر» می‌آید.
\end{warn}

\section{تغییر دادن عضوها}

برای عوض کردن یک عضو، آرایه باید \fcode{متغیر} باشد. سپس با اندیس، به همان
عضو مقدار می‌دهیم:

\program{تغییر عضو آرایه}
\begin{salam}
تابع اصلی:
    متغیر نمرات := [۱۷, ۱۵, ۱۹, ۱۲, ۲۰]
    چاپ "پیش از اصلاح:", نمرات
    نمرات[۳] = ۱۴
    نمرات[۰] = نمرات[۰] + ۱
    چاپ "پس از اصلاح:", نمرات
تمام
\end{salam}

\begin{outputfa}
پیش از اصلاح: [17, 15, 19, 12, 20]
پس از اصلاح: [18, 15, 19, 14, 20]
\end{outputfa}

هر عضو آرایه مثل یک متغیر معمولی رفتار می‌کند: می‌توان خواندش، در محاسبه
گذاشتش و به آن مقدار داد. تعداد عضوها اما ثابت است؛ آرایه‌ای که با پنج عضو
ساخته شده، تا آخر پنج عضو دارد.

\section{پیمایش آرایه با هر}

قدرت واقعی آرایه وقتی پیداست که با حلقه ترکیب شود. سلام حلقه‌ای مخصوص
آرایه دارد: \kwd{هر}. این حلقه، به تعداد عضوهای آرایه دور می‌زند و در هر
دور یک عضو را در متغیری می‌گذارد:

\program{پیمایش با هر}
\begin{salam}
تابع اصلی:
    میوه‌ها := ["سیب", "موز", "پرتقال"]
    هر میوه در میوه‌ها:
        چاپ "-", میوه
    تمام
تمام
\end{salam}

\begin{outputfa}
- سیب
- موز
- پرتقال
\end{outputfa}

\fcode{هر میوه در میوه‌ها:} را این‌طور بخوانید: «برای هر میوه در میوه‌ها».
در دور اول \fcode{میوه} برابر «سیب» است، در دور دوم «موز» و در دور سوم
«پرتقال». نه اندیسی لازم است، نه شمردن؛ سلام خودش همه را یکی‌یکی می‌دهد.
آرایه‌ی رشته‌ها هم درست مثل آرایه‌ی عددهاست؛ فقط همه‌ی عضوها باید هم‌نوع
باشند.

اگر شماره‌ی هر عضو را هم بخواهید، دو نام میان پرانتز بنویسید؛ اولی اندیس
را می‌گیرد و دومی عضو را:

\program{پیمایش با اندیس}
\begin{salam}
تابع اصلی:
    میوه‌ها := ["سیب", "موز", "پرتقال"]
    هر (شماره, میوه) در میوه‌ها:
        چاپ شماره + ۱, "-", میوه
    تمام
تمام
\end{salam}

\begin{outputfa}
1 - سیب
2 - موز
3 - پرتقال
\end{outputfa}

چون اندیس از صفر شروع می‌شود، برای نمایش شماره‌ی «انسانی» یکی به آن
اضافه کردیم.

\section{الگوهای پایه‌ی کار با آرایه}

بیشتر کارهایی که با آرایه می‌کنیم، به چند الگوی ساده برمی‌گردند. همه‌ی
آن‌ها از یک حلقه‌ی \fcode{هر} و یک یا دو متغیر تغییرپذیر ساخته می‌شوند.

\subsection{جمع و میانگین}

\program{جمع و میانگین نمره‌ها}
\begin{salam}
تابع اصلی:
    نمرات := [۱۷, ۱۵, ۱۹, ۱۲, ۲۰]
    متغیر مجموع := ۰
    هر نمره در نمرات:
        مجموع += نمره
    تمام
    تعداد := نمرات.طول()
    میانگین := (مجموع بعنوان اعشار۶۴) / (تعداد بعنوان اعشار۶۴)
    چاپ "مجموع:", مجموع
    چاپ "میانگین:", میانگین
تمام
\end{salam}

\begin{outputfa}
مجموع: 83
میانگین: 16.6
\end{outputfa}

همان الگوی جمع کردن فصل \ref{ch:loops} است، فقط این بار عددها از آرایه
می‌آیند. اگر فردا ده نمره‌ی دیگر به آرایه اضافه کنید، هیچ خط دیگری از
برنامه عوض نمی‌شود.

\subsection{بیشترین}

\program{گرم‌ترین روز هفته}
\begin{salam}
تابع اصلی:
    دماها := [۲۱, ۲۵, ۱۹, ۳۰, ۲۷, ۲۴, ۲۲]
    متغیر بیشترین := دماها[۰]
    متغیر روز بیشترین := ۰
    هر (روز, دما) در دماها:
        اگر دما > بیشترین:
            بیشترین = دما
            روز بیشترین = روز
        تمام
    تمام
    چاپ "گرم‌ترین روز: روز", روز بیشترین + ۱, "با دمای", بیشترین
تمام
\end{salam}

\begin{outputfa}
گرم‌ترین روز: روز 4 با دمای 30
\end{outputfa}

همان الگوی «بزرگ‌ترین سه عدد» فصل \ref{ch:conditions} است: با عضو اول
شروع می‌کنیم و هر عضوی که بزرگ‌تر بود، جایش را می‌گیرد. این بار اندیس عضو
برنده را هم نگه داشتیم تا بدانیم کدام روز بوده است. برای پیدا کردن کمترین
کافی است \code{>} را به \code{<} تبدیل کنید.

\subsection{شمارش}

\program{شمارش قبولی‌ها}
\begin{salam}
تابع اصلی:
    نمرات := [۱۷, ۹, ۱۹, ۱۲, ۸, ۲۰, ۱۴]
    متغیر قبولی := ۰
    هر نمره در نمرات:
        اگر نمره >= ۱۰: قبولی++ تمام
    تمام
    چاپ "قبول:", قبولی, "مردود:", نمرات.طول() - قبولی
تمام
\end{salam}

\begin{outputfa}
قبول: 5 مردود: 2
\end{outputfa}

\subsection{جست‌وجو}

آیا نام مشخصی در فهرست هست؟ برای این پرسش، یک متغیر منطقی می‌سازیم که
اول \fcode{نادرست} است و به محض پیدا شدن، \fcode{درست} می‌شود. وقتی پیدا
شد، با \fcode{بشکن} از ادامه‌ی جست‌وجو صرف نظر می‌کنیم:

\program{جست‌وجو در فهرست}
\begin{salam}
تابع اصلی:
    مهمانان := ["سارا", "رضا", "مینا", "علی"]
    جستجو := "مینا"
    متغیر پیدا شد := نادرست
    هر نام در مهمانان:
        اگر نام == جستجو:
            پیدا شد = درست
            بشکن
        تمام
    تمام
    اگر پیدا شد:
        چاپ جستجو, "در فهرست مهمانان هست"
    وگرنه:
        چاپ جستجو, "در فهرست مهمانان نیست"
    تمام
تمام
\end{salam}

\begin{outputfa}
مینا در فهرست مهمانان هست
\end{outputfa}

مقدار \fcode{جستجو} را به نامی که در فهرست نیست تغییر دهید و دوباره اجرا
کنید.

\section{آرایه به عنوان پارامتر تابع}

آرایه را می‌توان به تابع فرستاد. نوع پارامتر به شکل \fcode{صحیح[:]} نوشته
می‌شود که یعنی «آرایه‌ای از عددهای صحیح با هر طولی»، و هنگام فراخوانی،
بعد از نام آرایه \code{[:]} می‌نویسیم که یعنی «همه‌ی عضوها»:

\program{تابعی که آرایه می‌گیرد}
\begin{salam}
تابع مجموع(اعداد: صحیح[:]): صحیح:
    متغیر جمع := ۰
    هر عدد در اعداد:
        جمع += عدد
    تمام
    بازگشت جمع
تمام

تابع اصلی:
    کوچک := [۱, ۲, ۳]
    چاپ مجموع(کوچک[:])
    خریدها := [۴۵۰۰۰, ۱۲۰۰۰۰, ۳۲۰۰۰]
    چاپ "جمع خرید:", مجموع(خریدها[:])
تمام
\end{salam}

\begin{outputfa}
6
جمع خرید: 197000
\end{outputfa}

با این تابع، الگوی جمع کردن را یک بار برای همیشه نوشته‌ایم و برای هر
آرایه‌ای با هر طولی به کار می‌بریم. (نوشتن \code{[:]} با اعدادی درون
کروشه، مثل \fcode{خریدها[۱:۳]}، فقط بخشی از آرایه را می‌فرستد؛ در این
کتاب به آن نیازی نداریم.) تمرین‌های پایان فصل از شما می‌خواهند
همین کار را برای بیشترین و میانگین بکنید.

\section{چند آرایه‌ی هم‌راستا}

گاهی برای هر چیز چند اطلاعات داریم: نام دانش‌آموز و نمره‌اش، نام کالا و
قیمتش. راه ساده این است که برای هر اطلاعات یک آرایه بسازیم و قرارداد
کنیم که عضو شماره‌ی \fcode{ک} در همه‌ی آرایه‌ها به یک چیز مربوط است:

\program{نام‌ها و نمره‌ها}
\begin{salam}
تابع اصلی:
    نامها := ["سارا", "رضا", "مینا"]
    نمرات := [۱۸, ۱۲, ۱۵]
    هر (شماره, نام) در نامها:
        چاپ نام + ":", نمرات[شماره]
    تمام
تمام
\end{salam}

\begin{outputfa}
سارا: 18
رضا: 12
مینا: 15
\end{outputfa}

حلقه روی آرایه‌ی نام‌ها می‌چرخد و با اندیسی که می‌گیرد، نمره‌ی همان نفر را
از آرایه‌ی دوم برمی‌دارد. این روش ساده و کارآمد است، به شرطی که دو آرایه
همیشه هم‌طول و هم‌ترتیب بمانند.

\begin{deep}[وقتی تعداد از پیش معلوم نیست]
آرایه‌های این فصل طول ثابت دارند: نمی‌توان به آن‌ها عضو تازه افزود یا عضوی
را حذف کرد. برای فهرست‌هایی که در طول برنامه بزرگ و کوچک می‌شوند، سلام
ساختاری به نام \fcode{وکتور} دارد که با \fcode{بیفزا} عضو می‌گیرد و با
\fcode{دربیاور} عضو پس می‌دهد. همچنین برای نگه‌داری داده‌های مرتبط (مثلاً
نام و نمره و سن یک دانش‌آموز با هم) به جای چند آرایه‌ی هم‌راستا، می‌توان
\fcode{ساختار} تعریف کرد. این‌ها موضوع کتاب‌های بعدی سلام هستند؛ آنچه
در این فصل آموختید، پایه‌ی همه‌ی آن‌هاست.
\end{deep}

\section{یک برنامه‌ی کامل‌تر}

کارنامه‌ی یک کلاس کوچک: چاپ نمره و وضعیت هر نفر، میانگین کلاس و نام
بهترین دانش‌آموز.

\program{کارنامه‌ی کلاس}
\begin{salam}
تابع اصلی:
    نامها := ["سارا", "رضا", "مینا", "علی"]
    نمرات := [۱۸, ۱۲, ۱۵, ۱۹]
    متغیر مجموع := ۰
    متغیر بهترین := ۰

    چاپ "کارنامه‌ی کلاس"
    چاپ "----------------"
    هر (شماره, نام) در نامها:
        نمره := نمرات[شماره]
        وضعیت := نمره >= ۱۰ ؟ "قبول" : "مردود"
        چاپ نام + ":", نمره, "(" + وضعیت + ")"
        مجموع += نمره
        اگر نمره > نمرات[بهترین]: بهترین = شماره تمام
    تمام
    چاپ "----------------"
    تعداد := نمرات.طول()
    چاپ "میانگین کلاس:", (مجموع بعنوان اعشار۶۴) / (تعداد بعنوان اعشار۶۴)
    چاپ "بهترین دانش‌آموز:", نامها[بهترین]
تمام
\end{salam}

\begin{outputfa}
کارنامه‌ی کلاس
----------------
سارا: 18 (قبول)
رضا: 12 (قبول)
مینا: 15 (قبول)
علی: 19 (قبول)
----------------
میانگین کلاس: 16
بهترین دانش‌آموز: علی
\end{outputfa}

متغیر \fcode{بهترین} این بار \emph{اندیس} بهترین نمره را نگه می‌دارد نه
خود نمره را؛ به این ترتیب در پایان می‌توانیم با همان اندیس، نام او را از
آرایه‌ی نام‌ها بیرون بکشیم. مقایسه کنید با برنامه‌ی «گرم‌ترین روز» که هر
دو را نگه می‌داشت. هر دو راه درست است؛ انتخاب به این بستگی دارد که در
پایان به چه چیزی نیاز دارید.

\begin{summary}
\begin{itemize}
  \item آرایه فهرستی از مقدارهای هم‌نوع است: \fcode{[۱۷, ۱۵, ۱۹]}.
        \fcode{چاپ} کل آرایه را نشان می‌دهد و \fcode{.طول()} تعداد عضوها
        را.
  \item اندیس از صفر شروع می‌شود: \fcode{آ[۰]} اولین عضو و
        \fcode{آ[آ.طول() - ۱]} آخرین عضو است. اندیس بیرون محدوده خطای
        اجراست.
  \item برای تغییر عضوها آرایه باید \fcode{متغیر} باشد: \fcode{آ[۲] = ۵}.
  \item \fcode{هر عضو در آرایه:} همه‌ی عضوها را یکی‌یکی می‌دهد و
        \fcode{هر (اندیس, عضو) در آرایه:} اندیس را هم.
  \item الگوهای جمع، میانگین، بیشترین، شمارش و جست‌وجو همه از یک حلقه و
        یک متغیر تغییرپذیر ساخته می‌شوند.
  \item پارامتر آرایه‌ای با \fcode{صحیح[:]} نوشته می‌شود و آرایه با
        \fcode{نام[:]} به تابع فرستاده می‌شود.
\end{itemize}
\end{summary}

\section*{تمرین‌ها}

\exercise آرایه‌ای از قیمت چند کالا بسازید و گران‌ترین و ارزان‌ترین را
همراه با اختلافشان چاپ کنید.

\exercise تابعی به نام \fcode{میانگین} بنویسید که آرایه‌ای از عددهای صحیح
بگیرد و میانگین اعشاری آن‌ها را برگرداند. سپس تابعی به نام
\fcode{بیشترین} با همان الگو بنویسید.

\exercise آرایه‌ای از دمای هفت روز هفته بسازید و چاپ کنید چند روز دما از
میانگین هفته بالاتر بوده است. (راهنمایی: به دو حلقه نیاز دارید؛ اول
میانگین، بعد شمارش.)

\exercise آرایه‌ای از نام‌ها و آرایه‌ای هم‌راستا از سن‌ها بسازید. نام همه‌ی
کسانی را که زیر ۱۸ سال دارند چاپ کنید و در پایان بگویید مسن‌ترین نفر
کیست.

\exercise آرایه‌ای از ۱۰ عدد بسازید و یک آرایه‌ی تازه‌ی \fcode{متغیر} با
۱۰ صفر. با حلقه، آرایه‌ی دوم را وارونه‌ی آرایه‌ی اول پر کنید (عضو آخر اولی
در عضو اول دومی، و همین‌طور). هر دو را چاپ کنید.

\exercise برنامه‌ی زیر چه چاپ می‌کند؟ اول حدس بزنید، بعد اجرا کنید:
\begin{salam}
تابع اصلی:
    متغیر آ := [۱, ۲, ۳, ۴, ۵]
    هر (ک, عدد) در آ:
        اگر ک % ۲ == ۰: آ[ک] = عدد * ۱۰ تمام
    تمام
    چاپ آ
تمام
\end{salam}
